\par\noindent\SourceLeaf{51}%
\hspace*{\parindent}6. We can now proceed without delay to the proposed objective. Expand $(1+n\cos\varphi)^{\nu}$ by the binomial theorem into a series proceeding by powers of $\cos\varphi$, then resolve each term of this series into a new series proceeding by cosines of multiples of the given angle. Select from these series the terms multiplied by the cosine of the same angle and combine them into a single term. Once the law of the resulting terms is known, the whole task will be virtually complete. In the fundamental series

\[
\begin{aligned}
(1+n\cos\varphi)^{\nu} & = 1 + {}^{\nu}\mathfrak{B}^{(1} n\cos\varphi + {}^{\nu}\mathfrak{B}^{(2)} n^{2}\cos\varphi^{2} \\&\quad + {}^{\nu}\mathfrak{B}^{(3)} n^{3}\cos\varphi^{3} + .\,. + {}^{\nu}\mathfrak{B}^{(q)} n^{q}\cos\varphi^{q} + .\,.
\end{aligned}
\]

\noindent all the powers of $\cos\varphi$ occur in order, even and odd. When expanded, the even powers yield only cosines of even multiples and the odd powers only cosines of odd multiples. Thus we must turn to the latter if we are concerned with the cosine of an odd multiple, and to the former if we are concerned with the cosine of an even multiple.

A) The initial term of the series sought is obtained by taking, from all the proposed terms whose expansions begin with $\cos 0\varphi = 1$, their initial terms and combining them in their original order. Only even powers of $\cos\varphi$ yield such an initial term on expansion. Each contributes a part that belongs here. Thus part $h^{\text{ta}}$ comes from term $2h^{\text{to}}$ of the fundamental series, which is ${}^{\nu}\mathfrak{B}^{(2h)} n^{2h}\cos\varphi^{2h}$. Now for $\cos\varphi^{2h}$ the initial term, the only one that belongs here, is

\[
= \frac{1}{2^{2h-1}} \frac{{}^{2h}\mathfrak{B}^{(h)}}{2} = \frac{1}{2^{2h}} {}^{2h}\mathfrak{B}^{(h)};
\]

\noindent and so, on including the coefficient of $\cos\varphi^{2h}$, we obtain part $h^{\text{ta}}$ of the series multiplied by $\cos 0\varphi = 1$

\[
\frac{1}{2^{2h}} {}^{2h}\mathfrak{B}^{(h)}  .  {}^{\nu}\mathfrak{B}^{(2h)} n^{2h}.
\]

\noindent We have thus found the general term $h^{\text{tus}}$ of the series that constitutes the initial term of the series obtained by expanding $(1+n\cos\varphi)^{\nu}$ according to cosines of multiples of the angle. It is evident that its initial term is $=1$. If we put

\[
(1+n\cos\varphi)^{\nu} = A + A^{(1)}\cos\varphi + A^{(2)}\cos 2\varphi + .\,.\,.,
\]

\noindent then we have

\[
A = 1 + \frac{1}{2^{2}} {}^{2}\mathfrak{B}^{(1)}  .  {}^{\nu}\mathfrak{B}^{(2)} n^{2} + \frac{1}{2^{4}} {}^{4}\mathfrak{B}^{(2)}  .  {}^{\nu}\mathfrak{B}^{(4)} n^{4} + .\,.\,. + \frac{1}{2^{2h}} {}^{2h}\mathfrak{B}^{(h)}  .  {}^{\nu}\mathfrak{B}^{(2h)} n^{2h} + .\,.\,.
\]

\noindent or

\[
A = 1 + \frac{1}{2^{2}}  .  \frac{2}{1} \frac{\nu.\nu-1}{1.2} n^{2} + \frac{1}{2^{4}}  .  \frac{4.3}{1.2}  .  \frac{\nu.\nu-1.\nu-2.\nu-3}{1.2.3.4} n^{4} + \text{etc.}
\]
